\section{Mô hình đề xuất}

\subsection{Rời rạc hóa miền tính toán}
Mô hình đề xuất áp dụng trên một vùng không gian khối trụ rỗng (hollow cylinder) hoặc một phân đoạn hình trụ (cylindrical sector) trong hệ tọa độ trụ $(r, \theta, z)$, trong đó cấu trúc lưới được thiết lập để bao phủ hoàn toàn không gian tính toán của máy điện như minh họa trong Hình~\ref{fig1_a}. Miền tính toán $\mathbf{\Omega}$ được rời rạc hóa thông qua việc thiết lập ba tập hợp điểm nút lưới (grid points) thể hiện chi tiết tại Hình~\ref{fig1_b}, bao gồm: vector bán kính $\mathbf{R} = [r_0, r_1, \dots, r_{m}]$ xác định các lớp từ trục đến vỏ máy; vector vị trí góc $\mathbf{\Theta} = [\theta_0, \theta_1, \dots, \theta_{n}]$ rời rạc hóa theo phương tiếp tuyến; và vector trục $\mathbf{Z} = [z_0, z_1, \dots, z_{p}]$ rời rạc hóa không gian theo hướng trục.

\begin{figure}[htbp]
    \centering
    \begin{subfigure}[b]{0.45\textwidth}
        \centering
        \includegraphics[width=\textwidth]{figure/fig1a.png}
        \caption{Cấu trúc 3D tổng thể của động cơ AFPM và miền tính toán phân đoạn hình trụ.}
        \label{fig1_a}
    \end{subfigure}
    \vskip 0.5em
    \begin{subfigure}[b]{0.45\textwidth}
        \centering
        \includegraphics[width=\textwidth]{figure/fig1b.png}
        \caption{Rời rạc hóa phân đoạn hình trụ theo các trục $(r, \theta, z)$ và định danh phần tử $E_{i,j,k}$.}
        \label{fig1_b}
    \end{subfigure}
    \caption{Rời rạc hóa miền tính toán trong không gian ba chiều trên hệ tọa độ trụ.}
    \label{fig1}
\end{figure}

Dựa trên cấu trúc lưới này, toàn bộ miền tính toán được chia thành một tập hợp hữu hạn các phần tử hình học với tổng số lượng $q_{E}$ được xác định bởi:
\begin{equation}
    q_{E} = m \times n \times p
\end{equation}

Mỗi phần tử đơn lẻ $E_{i,j,k}$ (còn được ký hiệu ngắn gọn là $E_f$) được định danh duy nhất thông qua bộ chỉ số không gian $(i, j, k)$ với $0 \le i < m$, $0 \le j < n$ và $0 \le k < p$ như biểu diễn trong Hình~\ref{fig1_b}. Để thuận tiện cho việc ánh xạ dữ liệu vào hệ phương trình ma trận và tối ưu hóa tốc độ xử lý trong mã nguồn, mỗi phần tử còn được gán một chỉ số phẳng (flat index) $f$ duy nhất trong phạm vi $[0, q_E - 1]$. Với quy luật chỉ số $i$ tăng nhanh nhất, kế đến là $j$ và cuối cùng là $k$, chỉ số $f$ được xác định cụ thể theo quy tắc:
\begin{equation}
    f = k \cdot (n \cdot m) + j \cdot m + i
\end{equation}

Việc thiết lập mối quan hệ này giúp bộ giải 3D-MBGRN có thể chuyển đổi linh hoạt giữa cấu trúc hình học vật lý và các vector đại số trong quá trình tính toán.

Trong mô hình đề xuất, ma trận vị trí góc $\mathbf{\Theta}$ được thiết lập với bước góc đồng nhất $\Delta\theta = \text{const}$, tạo tiền đề toán học cho kỹ thuật dịch chỉ số (index-shifting) để mô tả chuyển động quay của rotor mà không cần tái tạo lưới.

\subsection{Cấu trúc phần tử lục diện 3D}

\subsubsection{Không gian giới hạn phần tử}
Trong mô hình 3D-MBGRN đề xuất, mỗi phần tử lục diện $E_{i,j,k}$ được giới hạn bởi sáu mặt đẳng tọa độ trong hệ tọa độ trụ $(r, \theta, z)$, bao gồm hai mặt trụ tại bán kính $r_i, r_{i+1}$, hai mặt phẳng kinh tuyến tại góc $\theta_j, \theta_{j+1}$, và hai mặt phẳng ngang tại cao độ $z_k, z_{k+1}$. Các phần tử liền kề tiếp xúc với nhau thông qua các mặt ranh giới chung mà không có sự giao cắt hay chồng lấn thể tích, cho phép mỗi phần tử được định danh một thuộc tính vật liệu độc lập và đồng nhất dựa trên vị trí tâm hình học. Đồng thời, cấu trúc này tạo điều kiện thuận lợi để đặt một nút thế từ vô hướng $P_{i,j,k}$ tại tâm của từng ô lưới phục vụ quá trình thiết lập hệ phương trình đại số toàn cục.

Dựa trên các ranh giới tọa độ đã thiết lập, các kích thước hình học đặc trưng của phần tử $E_{i,j,k}$ phục vụ cho việc tính toán mạch từ được xác định thông qua các độ dài và diện tích mặt cắt tương đương. Các chiều dài tương đương theo phương hướng tâm ($l_r$), tiếp tuyến ($l_\theta$) và dọc trục ($l_z$) được xác định chính xác theo các công thức sau:
\begin{equation}
    l_r = (r_{i+1} - r_i) \cos\left( \frac{|\theta_{j+1} - \theta_j|}{2} \right)
\end{equation}
\begin{equation}
    l_\theta = (r_{i+1} + r_i) \sin\left( \frac{|\theta_{j+1} - \theta_j|}{2} \right)
\end{equation}
\begin{equation}
    l_z = z_{k+1} - z_k
\end{equation}

Tương ứng với các độ dài vừa thiết lập, diện tích các mặt tương đương phục vụ cho quá trình bảo toàn thông lượng và tính toán khả năng dẫn từ của phần tử được xác định bởi:
\begin{equation}
    A_r = l_\theta \cdot l_z, \quad A_\theta = l_r \cdot l_z, \quad A_z = l_r \cdot l_\theta
\end{equation}

\subsubsection{Cấu tạo nội tại phần tử}
Mỗi phần tử $E_{i,j,k}$ được đặc trưng bởi một nút từ thế vô hướng đặt tại tâm hình học của khối lục diện ($P_{i,j,k}$) và sáu nhánh nội tại tỏa ra theo các trục tọa độ $s \in \{r, \theta, z\}$ như minh họa trong Hình~\ref{fig2_a}. Để quản lý dữ liệu mạng lưới, một quy ước chỉ số hướng $d \in \{0, 1\}$ được thiết lập đồng nhất cho từ trở ($R$), thông lượng từ ($\phi$) và nguồn sức từ động ($F$). Cụ thể, chỉ số $d=0$ đại diện cho các thông số trên nhánh đi vào nút thế trung tâm, trong khi $d=1$ đại diện cho các nhánh đi ra khỏi nút thế. Theo đó, các đại lượng nội tại trên mỗi nhánh được ký hiệu là $R_{s,d,i,j,k}$, $\phi_{s,d,i,j,k}$ và $F_{s,d,i,j,k}$.

\begin{figure}[htbp]
    \centering
    \begin{subfigure}[b]{0.45\textwidth}
        \centering
        \includegraphics[width=\textwidth]{figure/fig2a.png}
        \caption{Cấu tạo nội tại và các nhánh từ trở của phần tử lục diện đơn lẻ.}
        \label{fig2_a}
    \end{subfigure}
    \vskip 0.5em
    \begin{subfigure}[b]{0.45\textwidth}
        \centering
        \includegraphics[width=\textwidth]{figure/fig2b.png}
        \caption{Mạng lưới kết nối tương hỗ giữa các phần tử trong không gian 3D.}
        \label{fig2_b}
    \end{subfigure}
    \caption{Mô hình hóa mạng từ trở phần tử: (a) Cấu tạo nội tại mỗi phần tử và (b) Mạng lưới kết nối các phần tử.}
    \label{fig2}
\end{figure}

Về mặt vật lý, mỗi nhánh nội tại được mô hình hóa bằng một từ trở $R$ nối tiếp với nguồn sức từ động $F$. Khi các phần tử liền kề ghép nối với nhau trong không gian như thể hiện tại Hình~\ref{fig2_b}, các nhánh tiếp giáp qua mặt ranh giới chung sẽ liên kết trực tiếp với nhau, cho phép từ thông lưu chuyển liên tục giữa các phần tử trong toàn bộ miền tính toán.

Toàn bộ các phần tử trong không gian tính toán đều có cấu tạo và quy ước nhánh đồng nhất theo cùng một cấu trúc cơ sở. Tính đồng nhất này cho phép tự động hóa hoàn toàn quy trình sinh lưới và thiết lập các thông số phần tử một cách tổng quát, loại bỏ sự phụ thuộc vào các giả thiết mạch từ thủ công phức tạp.

\subsubsection{Xác định vật liệu và thuật toán gán tự động}
Trong mô hình 3D-MBGRN, mỗi phần tử $E_{i,j,k}$ được xem là một khối đồng nhất mang duy nhất một loại vật liệu. Thuộc tính vật liệu của phần tử được xác định thông qua tọa độ tâm hình học $(r_c, \theta_c, z_c)$ của nút thế $P_{i,j,k}$:
\begin{equation}
    r_c = \frac{r_i + r_{i+1}}{2}, \quad \theta_c = \frac{\theta_j + \theta_{j+1}}{2}, \quad z_c = \frac{z_k + z_{k+1}}{2}
\end{equation}

Thuật toán gán vật liệu tự động sử dụng quy trình kiểm tra vị trí điểm (point-in-region) so với đường bao hình học gốc của máy điện như minh họa trong Hình~\ref{fig456}. Nếu điểm tâm $(r_c, \theta_c, z_c)$ nằm trong phạm vi một miền vật lý xác định (như lõi sắt, nam châm hay khe hở không khí), toàn bộ phần tử $E_{i,j,k}$ sẽ được gán định danh vật liệu tương ứng.

\begin{figure}[htbp]
    \centering
    \includegraphics[width=0.45\textwidth]{figure/discrete3.png}
    \caption{Minh họa thuật toán gán vật liệu tự động theo vị trí điểm (point-in-region): Ánh xạ từ đường bao hình học thực tế (Geometry outlines) sang các ô lưới phần tử (Elements outlines) cho các vùng lõi sắt (Iron), nam châm (Magnet) và không khí (Air).}
    \label{fig456}
\end{figure}

Việc gán vật liệu đơn nhất theo điểm tâm có thể tạo ra hiệu ứng bậc thang (staircase effect) tại các bề mặt cong hoặc nghiêng phức tạp. Tuy nhiên, như được minh họa trong Hình~\ref{fig45}, sai số hình học này giảm dần và bề mặt vật liệu tiệm cận chính xác với mô hình gốc khi tăng mật độ chia lưới $(m, n, p)$.

\begin{figure}[htbp]
    \centering
    \includegraphics[width=0.48\textwidth]{figure/discrete1.png}
    \caption{Rời rạc hóa hình học máy điện AFPM theo các cấp độ lưới khác nhau: So sánh giữa mô hình hình học gốc (Original) và các cấu hình lưới 3D-MBGRN với mật độ tăng dần từ thô đến mịn (5.460 phần tử, 24.480 phần tử và 99.360 phần tử).}
    \label{fig45}
\end{figure}

\subsubsection{Xác định thông số từ trở cho các nhánh nội tại}
Từ kích thước hình học tương đương và thuộc tính vật liệu đã gán, từ trở của sáu nhánh nội tại bên trong phần tử $E_{i,j,k}$ được xác định từ nút thế trung tâm $P_{i,j,k}$ đến tâm sáu mặt ranh giới. Do nút thế đặt tại tâm hình học, chiều dài đường từ của mỗi nhánh bằng một nửa chiều dài tương đương $l_s$ theo phương tương ứng ($s \in \{r, \theta, z\}$):
\begin{equation}
    R_{s,d,i,j,k} = \frac{l_s}{2 \mu_0 \mu_{r,i,j,k} A_s}, \quad s \in \{r, \theta, z\}, \; d \in \{0, 1\}
\end{equation}
trong đó $\mu_0$ là độ từ thẩm chân không, $A_s$ là diện tích mặt cắt ngang tương đương vuông góc với phương $s$, và $\mu_{r,i,j,k}$ là độ từ thẩm tương đối của phần tử.

Đối với vật liệu sắt từ phi tuyến, $\mu_{r,i,j,k}$ phụ thuộc vào mật độ từ thông $B_{s,d,i,j,k}$ trên từng nhánh và được tính toán thông qua từ thông $\phi_{s,d,i,j,k}$:
\begin{equation}
    B_{s,d,i,j,k} = \frac{\phi_{s,d,i,j,k}}{A_s}
\end{equation}
Giá trị $\mu_{r,i,j,k}$ sau đó được cập nhật lặp thông qua đường cong từ hóa $B\text{-}H$ của vật liệu trong quá trình giải hệ phương trình phi tuyến toàn cục.

\subsubsection{Nguồn kích thích từ tính}
Tại mỗi phần tử $E_{i,j,k}$, vector nguồn sức từ động (MMF) tổng hợp $\mathbf{F}_{i,j,k} = [F_{r,i,j,k}, F_{\theta,i,j,k}, F_{z,i,j,k}]$ được xác định bởi sự đóng góp từ thành phần nam châm vĩnh cửu ($\mathbf{F}_{m,i,j,k}$) và nguồn kích thích từ hệ thống dây quấn ($\mathbf{F}_{w,i,j,k}$):
\begin{equation}
    \mathbf{F}_{i,j,k} = \mathbf{F}_{m,i,j,k} + \mathbf{F}_{w,i,j,k}
\end{equation}
Trong đó, $\mathbf{F}_{i,j,k}$ đại diện cho tổng tác động kích từ cục bộ tại vị trí phần tử, đóng vai trò là nguồn sơ cấp gây ra sự chênh lệch thế từ nút trong mạng lưới.

Nguồn sức từ động do nam châm vĩnh cửu bên trong phần tử $E_{i,j,k}$ được xác định bởi:
\begin{equation}
    \mathbf{F}_{m,i,j,k} = (H_{c,i,j,k} \cdot l_{m,i,j,k}) \cdot \mathbf{u}_{m,i,j,k}
\end{equation}
trong đó $H_{c,i,j,k}$ là cường độ lực kháng từ, $l_{m,i,j,k}$ là chiều dài tương đương của phần tử dọc theo phương từ hóa, và $\mathbf{u}_{m,i,j,k}$ là vector đơn vị chỉ hướng từ hóa.

Đối với dây quấn, sức từ động cục bộ bên trong phần tử $E_{i,j,k}$ được xác định bởi:
\begin{equation}
    \mathbf{F}_{w,i,j,k} = (\mathbf{N}_{w,i,j,k} \odot \mathbf{I}_{i,j,k}) \cdot \mathbf{n}_w
\end{equation}
trong đó $\mathbf{N}_{w,i,j,k} = [n_1, n_2, \dots, n_m]$ và $\mathbf{I}_{i,j,k} = [i_1, i_2, \dots, i_m]$ lần lượt là vector số vòng dây và vector dòng điện tức thời của $m$ pha; $\odot$ ký hiệu phép nhân Hadamard biểu thị tổng ampere-vòng; và $\mathbf{n}_w = [\mathbf{n}_{w,1}, \mathbf{n}_{w,2}, \dots, \mathbf{n}_{w,m}]^T$ là vector pháp tuyến đơn vị chỉ hướng trục của các cuộn dây.

Tại mỗi phần tử $E_{i,j,k}$, sức từ động các phương $r$, $\theta$, $z$ được phân bổ đối xứng trên các nhánh tương ứng:
\begin{equation}
    F_{s,0,i,j,k} = F_{s,1,i,j,k} = \frac{\mathbf{F}_{s,i,j,k}}{2}, \quad s \in \{r, \theta, z\}
\end{equation}

\subsection{Điều kiện biên}
Để đóng kín hệ phương trình đại số và đảm bảo tính duy nhất của nghiệm cho mô hình 3D-MBGRN, các điều kiện biên phải được thiết lập dựa trên đặc tính vật lý và tính đối xứng hình học của máy điện. Việc áp dụng các điều kiện biên phù hợp không chỉ giúp phản ánh chính xác trường điện từ mà còn cho phép giảm đáng kể kích thước ma trận hệ thống, từ đó cải thiện tốc độ tính toán.

Điều kiện biên song song (parallel boundary condition) được áp dụng một cách tự nhiên tại các mặt giới hạn theo phương hướng tâm ($r$) và phương dọc trục ($z$). Cụ thể, đối với các phần tử nằm tại ranh giới trong ($i=0$), ranh giới ngoài ($i=m$), hoặc các mặt cắt dọc trục ($k=0$ và $k=p$), thông lượng từ $\phi_{s,d,i,j,k}$ xuyên qua các mặt ranh giới này được cưỡng bức bằng không:
\begin{equation}
\phi_{s, d, i, j, k} = 0 \quad \text{khi} \quad
\begin{cases} 
s = r, & (d, i) \in \{(0, 0), (1, m)\} \\ 
s = z, & (d, k) \in \{(0, 0), (1, p)\} 
\end{cases}
\end{equation}

Theo phương tiếp tuyến ($\theta$), các điều kiện biên đối xứng được áp dụng để thu hẹp miền tính toán. Miền đối xứng chu kỳ cơ bản của máy điện được xác định bởi ước chung lớn nhất giữa số rãnh stator ($Q_s$) và số cực từ ($2p$), tương ứng với góc mở $\theta_{\text{per}} = \frac{2\pi}{\gcd(Q_s, 2p)}$~\cite{zhu2002improved}. Nếu số cực trong một đơn vị cơ bản $2p' = \frac{2p}{\gcd(Q_s, 2p)}$ là một số lẻ và cấu trúc phân bố dây quấn thỏa mãn tính chất đối xứng ngược pha, điều kiện biên phản chu kỳ (anti-periodic) có thể được áp dụng để giảm quy mô miền tính toán xuống còn một nửa. Do đó, độ mở góc tối thiểu $\theta_{\min}$ của mô hình 3D-MBGRN được xác định cụ thể theo:
\begin{equation}
    \theta_{\min} = 
    \begin{cases} 
        \dfrac{\pi}{\gcd(Q_s, 2p)}, & \text{với điều kiện biên phản chu kỳ} \\[8pt]
        \dfrac{2\pi}{\gcd(Q_s, 2p)}, & \text{với điều kiện biên chu kỳ}
    \end{cases}
\end{equation}
Việc giới hạn không gian khảo sát trong phạm vi $\theta_{\min}$ giúp giảm đáng kể số lượng biến thế nút cần giải trong hệ phương trình toàn cục mà vẫn bảo toàn đầy đủ các đặc tính trường điện từ của máy điện.

Tùy thuộc vào đặc điểm đối xứng, kết nối tại biên đầu ($j=0$) và biên cuối ($j=n$) theo phương $\theta$ được thiết lập. Với điều kiện biên chu kỳ, thông lượng được bảo toàn hoàn toàn về độ lớn và chiều:
\begin{equation}
    \phi_{\theta,0,i,0,k} = \phi_{\theta,1,i,n,k}
\end{equation}
Ngược lại, điều kiện biên phản chu kỳ được áp dụng khi từ trường có tính chất đảo cực tại biên đối xứng, dẫn đến sự đảo dấu của thông lượng:
\begin{equation}
    \phi_{\theta,0,i,0,k} = -\phi_{\theta,1,i,n,k}
\end{equation}

\begin{figure}[htbp]
    \centering
    \includegraphics[width=0.48\textwidth]{figure/boundary.png}
    \caption{Minh họa cách gán điều kiện biên}
    \label{fig4524}
\end{figure}

Thông qua việc thiết lập các điều kiện biên chu kỳ và phản chu kỳ, từ thông được phép lưu chuyển trực tiếp và liên tục từ mặt phẳng ranh giới $\theta$ này sang mặt phẳng ranh giới $\theta$ đối diện, tạo nên một hệ thống mạng từ trở khép kín hoàn chỉnh. 

\subsection{Mô phỏng chuyển động quay}
Trong mô hình 3D-MBGRN trên hệ tọa độ trụ, miền không gian được chia với bước góc đồng nhất $\Delta\theta = \theta_{j+1} - \theta_j$. Do đó, toàn bộ các phần tử $E_{i,j,k}$ trên cùng một bán kính $r_i$ và cao độ $z_k$ đều có kích thước hình học giống hệt nhau theo phương tiếp tuyến.

Đặc tính đồng nhất này cho phép mô phỏng chuyển động quay của rotor qua từng bước góc $s \cdot \Delta\theta$ ($s \in \mathbb{Z}$) thông qua phép dịch chỉ số phần tử:
\begin{equation}
    E_{i,j,k} \leftarrow E_{i,(j+s) \pmod n, k}
\end{equation}

\begin{figure}[htbp]
    \centering
    \includegraphics[width=0.48\textwidth]{figure/rotate.png}
    \caption{Minh họa kỹ thuật dịch chỉ số lưới (index-shifting) để mô phỏng chuyển động quay trên máy điện từ thông hướng kính (trái) và máy điện từ thông dọc trục (phải).}
    \label{fig452}
\end{figure}

Như minh họa trong Hình~\ref{fig452}, chuyển động quay được thực hiện bằng cách hoán vị thuộc tính vật liệu của các phần tử rotor mà không làm thay đổi cấu trúc liên kết tĩnh của mạng từ trở. Kỹ thuật dịch chỉ số loại bỏ hoàn toàn nhu cầu tái tạo lưới (re-meshing) hoặc sử dụng các phép nội suy tại mặt tiếp xúc trượt, qua đó giảm thiểu đáng kể chi phí xử lý và nâng cao hiệu năng tính toán cho bộ giải.